\documentclass[11pt,letterpaper]{article}
\PassOptionsToPackage{hyphens}{url}
\usepackage[margin=.83in,top=.8in,bottom=.8in,headheight=15pt]{geometry}
\usepackage[T1]{fontenc}
\usepackage{lmodern,amsmath,amssymb,mathtools,bm,booktabs,array,graphicx,tikz}
\usetikzlibrary{arrows.meta,positioning,calc,decorations.pathreplacing,patterns}
\definecolor{ink}{HTML}{17344D}\definecolor{accent}{HTML}{087E8B}\definecolor{pale}{HTML}{EDF5F5}\definecolor{gold}{HTML}{AD6B19}
\providecommand{\dd}{\,\mathrm{d}}\providecommand{\R}{\mathbb{R}}
\providecommand{\E}{\mathbb{E}}\providecommand{\Var}{\operatorname{Var}}
\providecommand{\Cov}{\operatorname{Cov}}\providecommand{\diag}{\operatorname{diag}}
\providecommand{\norm}[1]{\left\lVert#1\right\rVert}
\providecommand{\abs}[1]{\left|#1\right|}
\usepackage{microtype,hyperref,fancyhdr,enumitem,amsthm,sectsty,needspace}
\allsectionsfont{\raggedright}
\hypersetup{colorlinks=true,linkcolor=ink,urlcolor=accent,pdfauthor={Zan Ahmad},pdftitle={Least squares, principal components, and clustering / Review}}
\setlength{\parindent}{0pt}\setlength{\parskip}{6pt}\linespread{1.07}
\setlist{itemsep=3pt,topsep=4pt,leftmargin=*}
\emergencystretch=2em\clubpenalty=10000\widowpenalty=10000
\pagestyle{fancy}\fancyhf{}\fancyhead[L]{\small Scientific Computing in Python}\fancyhead[R]{\small Lecture 5}
\fancyfoot[L]{\footnotesize Zan Ahmad}\fancyfoot[R]{\footnotesize\thepage}\renewcommand{\headrulewidth}{.3pt}
\newtheorem{proposition}{Proposition}\newtheorem{theorem}{Theorem}\newtheorem{lemma}{Lemma}
\theoremstyle{definition}\newtheorem{definition}{Definition}\newtheorem{example}{Example}
\numberwithin{equation}{section}
\begin{document}
\section*{Least squares, principal components, and clustering --- review sheet}
Use this after attempting the worked examples and practice. Explain each claim in words before calculating.

\section*{Key ideas}
\begin{itemize}
\item Choose a scientific output before selecting a method.
\item Regression uses a target; PCA and clustering here do not.
\item Report test error with its units, split and simple baseline.
\item Centering and scaling are fitted transformations.
\item Projection can discard information, while a full basis change need not.
\item Cluster labels are arbitrary; metric and number of centers are choices.
\end{itemize}
\section*{Common errors and repairs}
\textbf{Fitting preprocessing on all rows before a test split.} Split first and fit transformations on training rows.

\textbf{Reading explained variance as prediction accuracy.} It describes feature variation under the chosen metric.

\textbf{Treating raw cluster numbers as class predictions.} Group identities require external meaning; label permutations are equivalent.

\textbf{Claiming a unit conversion changed the specimens.} It changed raw numerical geometry; state the metric or standardization.

\section*{Retrieve without your notes}
\begin{itemize}
\item What are the units of a response residual?
\item Why can a PCA direction change sign without changing reconstruction?
\item What does a rank-deficient line fit reveal?
\item What is needed to interpret a cluster scientifically?
\end{itemize}
\section*{Readiness check}
\begin{itemize}
\item Fit a straight-line baseline with an explicit response and coefficient units.
\item Calculate held-out RMSE and compare it with a training-mean baseline.
\item Center and project two-feature data while explaining what PCA preserves.
\item Trace one two-center assignment/update and distinguish cluster labels from known classes.
\item Fit transformations on training rows and test the effect of changing a feature unit.
\end{itemize}
\end{document}
